\clearpage
\chapter*{结\quad 论}
\markboth{结论}{结论}
\phantomsection
\addcontentsline{toc}{chapter}{结\quad 论}

本文围绕校园信息分散、重复咨询多、通用大模型回答缺乏校本依据等问题，设计并实现了一个基于大模型的校园智能问答系统。系统以学校公开的规章制度、办事指南和常见问答为知识源，把“资料入库—问题理解—混合召回—融合重排—受约束生成—引用返回”串成一条完整链路，力求让回答既有自然语言的易读性，又能追溯到具体资料。

在需求分析阶段，本文从技术、经济、操作和进度四个方面论证了可行性，划分了学生、教师、管理员和超级管理员四类角色，梳理了登录与权限、资料上传与分类、解析与知识库构建、自然语言问答、引用溯源、历史与反馈、统计与配置等功能需求，并提出了响应时间、并发、可靠性、安全性、可维护性和可扩展性等非功能要求。在系统设计阶段，本文采用前后端分离的分层架构，把表现层、接口层、业务层、数据与模型层分离，设计了结构感知分块、BM25 与向量双路召回、RRF 融合、交叉编码器重排、提示词约束、引用校验和分级拒答等关键流程，并给出了数据库表、RESTful 接口和权限安全方案。在实现阶段，本文说明了 FastAPI 认证依赖、文档解析与分块、RRF 融合、提示词组装和引用校验等核心代码结构，以及管理统计和配置的实现方式。

本文的主要工作可以概括为三点。第一，针对校园文档的标题层级和条款结构，设计了保留标题路径、页码和权限标签的分块方法，使检索片段既能保持语义完整，又能定位到原文。第二，把稀疏检索和稠密检索结合起来，先用 RRF 融合候选，再用交叉编码器重排，在不显著增加成本的前提下提高顶部证据的相关性。第三，把引用溯源和拒答作为生成流程的组成部分，通过提示词约束、引用编号校验和分级降级，降低模型编造内容的概率，让用户能够核对答案来源。

本文还存在一些不足。首先，系统核心模块的实现方案和测试方案已经形成，但完整的联调、性能测试和对照实验仍在推进，实测数据需要在后续工作中补充；本文没有用估计值代替实验结果。其次，知识库目前以文本类文档为主，对扫描版 PDF、图片表格和富媒体内容的支持不够，后续可以接入 OCR 和多模态模型。再次，嵌入模型和重排模型在 CPU 上运行，面对大规模知识库时检索时延会上升，需要通过模型量化、缓存或 GPU 推理优化。最后，多轮对话的问题改写仍依赖大模型，改写错误可能影响召回，后续可以结合规则和用户反馈做更稳健的处理。

后续工作可以从四个方向展开：一是扩充校园知识库，建立资料更新和责任分工机制，保证知识库与学校最新制度同步；二是完善测试集和标注规范，扩大样本规模，增加跨主题、对比式和否定式问题，提高评价的覆盖面；三是研究自适应检索，根据问题类型动态调整召回通道、候选数量和重排阈值，减少不必要的模型调用；四是把系统部署到校园内网，结合实际使用数据观察用户行为，持续改进回答质量和引用展示。通过上述工作，系统有望从毕业设计原型逐步发展为可实际使用的校园信息服务平台。
